\SourceLeaf{385}

\par\bigskip\Needspace{8\baselineskip}
\pdfbookmark[1]{Remarks on the Theory of Biquadratic Residues}{nachlass-biquadratic-bemerkungen}
\begin{center}\large\scshape Remarks.\end{center}
\BookRule{.2\textwidth}
Judging by the position of the manuscripts in a notebook, the fragments designated I and II in this printed edition date from $1811$ or the time immediately following. Of the preceding attempts to prove the fundamental theorem for biquadratic residues by the methods applied here to the residue $1+i$, a record survives that settles the special case of the residue $1+2i$ and starts from the determination of the biquadratic character appended as a note to Article $6$ of Fragment III. The manuscripts also supplement the historical information published by \textsc{Gauss} in the notices of his arithmetical papers: the theorems included in Articles $15$ to $20$ of the \emph{Theoria residuorum biquadrat.} were written down before the \emph{Theoria motus corporum coel.} was worked out. The investigations of cubic residues mentioned in those notices were probably never worked up in full; what is recorded consists of proofs of the reciprocity theorems for two primes, one of them real, carried out with the methods provided by the paper \emph{Disquisitionum circa aequationes puras ulterior evolutio}.

Fragments III to VI form separate booklets in the manuscripts. For the first three, the handwriting points to a date not far from that of Fragments I and II, while a considerably later date must be assumed for the last, No. VI.

\par\medskip
[I.] Art. $10$. The number of lattice points in $(z,z',z'+\zeta,z+\zeta)$ is determined by the theorem: if $a$ and $b$ are relatively prime, the straight line drawn from $\xi+\eta i$ to $\xi+\eta i+a+bi$ passes through one lattice point when the imaginary part of $(\xi+\eta i)\pdot(-a+bi)$ is an integer.

\SourceLeaf{386}

[I.] Art. $17$. The first transformation of the last $S$ in the expression for $\Delta S$ is obtained by dividing the region in question into three parts such that this $S$ becomes
\[
\begin{aligned}
&S\bigl(\tfrac12m+\tfrac12-\tfrac12i,\ (\tfrac12+\tfrac12i)m,\\
&\qquad(\tfrac12+\tfrac12i)m+i,\ \tfrac12m+\tfrac12+\tfrac12i_{(+)}\bigr)\\
&-S\bigl(\tfrac12m,\ \tfrac12m+\tfrac12+\tfrac12i_{(+)},\ \tfrac12m+\tfrac12-\tfrac12i\bigr)\\
&-S\bigl((\tfrac12+\tfrac12i)m,\ \tfrac12m,\ \tfrac12m+\tfrac12-\tfrac12i\bigr)
\end{aligned}
\]
then counting the lattice points in the first part by means of the theorem in Article $10$, and further noting that the second part contains no lattice point.

The second transformation is obtained by multiplying the quantities corresponding to the vertices of the third part by $i$ and increasing them by the integer $(1-i)\frac{m-1}{2}$; finally, the third is obtained by comparing the resulting figure, according to the rule of Article $16$, with the figure displaced from it by $\frac{-1}{1+i}$.

[I.] Art. $[18.]$ An explanation of the first scheme is found in Articles $1$ to $6$ of the fragment written later and designated here by [II.].

[I.] Art. $[18.]\ (2.)$ The geometric interpretation, with the aid of the two systems of lattice points
\[
[-2iQ-\tfrac12,\tfrac12,iQ]=[-2iQ+\tfrac12,\tfrac12,iQ]=\mathrm{IV}^{*}
\]
and
\[
-[-2iQ+\tfrac12i,-\tfrac12i,iQ]=-[-2iQ-\tfrac12i,-\tfrac12i,iQ]=\mathrm{X}^{*}
\]
gives the equations
\[
\begin{aligned}
\mathrm{IV}-\mathrm{IV}^{*}+[-2iQ,\tfrac12i,-\tfrac12]
&=-\mathrm{X}+\mathrm{X}^{*}+[-iQ,\tfrac12i,-\tfrac12]\\
\mathrm{XIII}+\mathrm{IV}^{*}&=[-2iQ,1,iQ]=-I(-2iQ)+I(-iQ)\\
\mathrm{V}+\mathrm{X}^{*}&=-[-2iQ,-i,iQ]=R(-2iQ)-R(-iQ)
\end{aligned}
\]
where, in general, $Rx$ and $Ix$ denote the greatest integers contained in the real part and in the coefficient of the imaginary part of $x$, respectively.

But $[-2iQ,\tfrac12i,-\tfrac12]$ is the number of lattice points in the square whose centre is at $-2iQ-\tfrac14+\tfrac14i$ and whose endpoints have the displacements $\tfrac12$ and $\tfrac12i$ between them.

$[-iQ,\tfrac12i,-\tfrac12]$ or $[-2iQ,-\tfrac12i,\tfrac12]$ is the number of lattice points in an equal square with centre $-2iQ+\tfrac14-\tfrac14i$.

[I.] Art. $[18.]\ (3.)$ With the aid of the lattice points $[0,\tfrac12i,-iQ]=-[-\tfrac12i,-\tfrac12i,-iQ]=\mathrm{I}^{*}$ one obtains

\SourceLeaf{387}

\[
\begin{aligned}
\mathrm{VII}-\mathrm{XII}&=\mathrm{I}-\mathrm{I}^{*}+[-iQ,\tfrac12,\tfrac12i]-[0,\tfrac12,\tfrac12i]\\
&=\mathrm{I}-\mathrm{I}^{*}+[-2iQ,\tfrac12,\tfrac12i]\\
\mathrm{II}-\mathrm{I}^{*}&=[0,-i,-iQ]=-R0+R(-iQ)\\
\mathrm{IX}-\mathrm{I}&=[0,-1,-iQ]=I0-I(-iQ)\\
\mathrm{VI}-\mathrm{XII}&=-[\tfrac12i,-1,-iQ]=-I(\tfrac12i)+I(-iQ+\tfrac12i)\\
\mathrm{VIII}-\mathrm{XIV}&=-[\tfrac12,-1-i,\tfrac12im_{(-)}]=\frac{a-1}{2}+\frac b2
\end{aligned}
\]

[II.] Art. $10$. We have
\[
\begin{aligned}
\Upsilon P'&\equiv-\Upsilon(-iP'+\tfrac12im),\\
\Upsilon P''&\equiv-1-\Upsilon(P''-\tfrac12im),\\
\Upsilon P'''&\equiv1+\Upsilon(-iP''')\ (\text{mod. }4)
\end{aligned}
\]
and $-iP'+\tfrac12im$, $P''-\tfrac12im$ are the lattice points in I and VI, respectively, increased by $\tfrac12i$.

[III.] Art. $6$. I have added to the present paper the rules given in the note for determining $\mathrm{Dec.}\frac Mm$, taking them from another place in the manuscripts. The first of these two rules, which, as is easy to see, agrees with the second, follows from that of Article $6$, because
\[
\begin{gathered}
k\equiv f\pdot i^{-n}\ (\text{mod. }m),\quad n=\Theta\frac{2fM}{m},\quad p=mm',\\
\left[\frac{2km'}p\right]^2\equiv\Theta\frac{2km'}p\ (\text{mod. }2+2i)
\end{gathered}
\]
holds.

[III.] Art. $8$ contains in the manuscript an example for Article $7$, namely the determination of the decident of $-1+2i$ for the modulus $-11+4i$.

[III.] Art. $10$. For the remark ‘to be expressed differently’, compare Article $3$ of the following fragment [IV].

[IV.] Articles $1$. $2$. $4$ contain in the manuscript, besides what is printed here, the application to the two examples for $m=5+8i$, $M=9+4i$ and for $m=9+4i$, $M=5+8i$.
